\chapter{Completeness of dynamic-priority dialogues}
\label{sec:p18-is-dynamic-completeness}

The dynamic semantics in Chapter~\ref{sec:p18-is-dynamic} permits
several compatible justified arguments to supply priorities.
The original dialogue, however, requires its next response to be
one argument whose own conclusions supply the comparison.
These requirements are not equivalent over the source's full
strict-rule domain. This chapter gives a finite counterexample
and a corrected dialogue with one argument per move.
The correction retains the original arguments and dynamic semantics,
but permits a suspended challenge to receive separately established
lemmas. Every rule change is stated below.

The historical target is Definitions 6.1--6.2 and
Propositions 6.3--6.5 of the Prakken--Sartor technical companion,
particularly its Lemma 10.3 and the proof of Proposition 6.4
on author pages 46--47
\cite{PrakkenSartor1997Priorities}.
The related 1996 discussion supplies the neutralization motivation
\cite{PrakkenSartor1996Dialectical}.
The earlier support-explicit certificates are not used as the
definition of the game in this chapter.

\section{The obstruction to a single self-supported response}

Retain \(\mathcal A\), \(C_S(A)\), attack, coherence,
\(E_X,D_X,D_X^s,G,J_\alpha,J\) from
Chapter~\ref{sec:p18-is-dynamic}. In particular, strict rules need
not be closed under contraposition. Every argument has finite
content and a unique producing rule for each of its conclusions.
The empty argument has no conclusions when the strict rules have
no empty-premise consequences.

\begin{hbdefinition}{def:p18-is-compression-countertheory}{A finite theory with incompatible amalgamations}
Take four distinct ordinary atoms \(a,b,c,d\) and eight distinct defeasible
rule names \(u,v,t,z,m,n,p,q\).
The ordinary signature contains exactly those four atoms.
The defeasible rules are exactly
\[
\begin{array}{llll}
u:\Rightarrow a,&v:\Rightarrow b,&t:\Rightarrow c,&z:\Rightarrow d,\\
m:\Rightarrow\neg c,&n:\Rightarrow\neg d,&
p:a\Rightarrow m\triangleleft t,&
q:b\Rightarrow n\triangleleft z.
\end{array}
\]
Besides all ground instances of the four strict priority rules,
take precisely the ordinary strict rules
\[
 a,b\to\neg a,\qquad c,a\to\neg c,\qquad d,b\to\neg d.
\]
All priority variables range over the finite set of displayed
defeasible names. There are no function symbols, other defeasible
rules, weak premises or empty-premise strict rules.
This is a finite grounding of the source language. The ordinary
strict rules contain no priority predicate and the strict base
alone is consistent.

Use the legal argument contents
\[
 U=[u],\quad V=[v],\quad P=[u,p],\quad Q=[v,q],
 \quad T=[t],\quad Z=[z],\quad A=[t,z],\quad B=[m,n].
\]
The capitals name arguments; the corresponding lower-case letters
name rules or ordinary atoms as specified above.
\end{hbdefinition}

\begin{hbtheorem}{thm:p18-is-compression-countermodel}{proved=annex:p18.is-compression-countermodel}{A justified argument without an original-game proof}
In this theory
\[
 [],V,Q\in J_1,\qquad U,P,Z\in J_2,\qquad T,A\in J_3.
\]
The support set \(\{P,Q\}\) is conflict-free, but every argument
concluding both of its positive priority conclusions is incoherent.
The argument \(B\) is coherent and \(B D_\varnothing A\).
Every original-game response \(C\) to this challenge, satisfying
\[
 C D_{\{C\}}^s B
 \quad\text{or}\quad
 \neg(B D_{\{C\}}A),
\]
is incoherent. The opponent can then play the empty argument,
against which no response is possible.
Thus the original single-response game is incomplete even for
finite theories, with global reply minimality removed and even
if repeated proponent arguments are allowed.
\end{hbtheorem}

The failure is not caused by an illegal concatenation alone.
For example, \(Z+Q=[z,v,q]\) is a legal argument, with distinct
heads and all positive premises supplied. It is nonetheless
incoherent: \(d,b\to\neg d\) supplies a strict contradiction.
Although \(Z\) and \(Q\) are individually justified, their combination
cannot be justified.

\begin{hbtheorem}{thm:p18-is-compression-rank}{proved=annex:p18.is-compression-rank}{The earlier-rank merger implication fails}
For the same \(B\), \(Z\) strictly defeats \(B\) relative to
\(\{Q\}\). Moreover \(B\) does not attack \(Q\), and
\(Z+Q\) strictly defeats \(B\) relative to its own conclusions.
Nevertheless \(Q\in J_1\), \(Z\in J_2\), and \(Z+Q\notin J\).
Consequently the absence of an attack on the added support does
not justify the inference that the merged response belongs to
the earlier stage. Finite priority information by itself does
not prove that inference.
\end{hbtheorem}

The example remains a counterexample if an additional default
\(u':\Rightarrow a\) is admitted. An argument for the first priority
may then use either producer of \(a\), but cannot include both
producers of that same head. The contradiction arguments still
apply to every legal choice. This variation separates repeated
heads from the more basic coherence failure.

\section{A corrected dialogue with sequential lemma hearings}

\begin{hbdefinition}{def:p18-is-licensing-state}{States, players, and licensed arguments}
There are two players, proponent \(P_{\rm pl}\) and opponent
\(O_{\rm pl}\). The subscript distinguishes a player from the
particular counterexample argument \(P=[u,p]\).
A state contains an ordered list \(L\) of distinct
\emph{licensed arguments} and a nonempty stack of frames.
Each frame records a legal claim argument \(C\) and is either
waiting for an opponent challenge or contains a particular open
challenge \(B\). Only the top frame is active.
The initial state has empty \(L\) and one frame for the target \(A_0\),
waiting for a challenge.

Write \(|L|\) for the set of entries of \(L\), and abbreviate
\(D_{|L|}\) by \(D_L\) and \(E_{|L|}\) by \(E_L\).
The list does not merge its entries into
one argument. It extracts their actual priority conclusions using
the existing definition of \(E_{|L|}\).
The adjective licensed describes the dialogue record; it does
not assume in advance that an entry is semantically justified.
That implication for winning play will be proved.
Move tests use the raw formula \(D_{|L|}\) on every finite list;
they do not silently assume that an arbitrary recorded list
is conflict-free.
\end{hbdefinition}

\begin{hbdefinition}{def:p18-is-licensing-moves}{Challenge, hearing, resolution, and return}
At a waiting frame with claim \(C\), the opponent chooses any
\(B D_\varnothing C\). That frame then records the open challenge
\((C,B)\). If no such \(B\) exists, the frame closes automatically.

At an open challenge \((C,B)\), the proponent has three kinds of
move, each naming one legal argument \(H\):
\begin{enumerate}
\item \emph{Hear \(H\)}: \(H\) is neither in \(L\) nor the claim
of any currently pending frame. Suspend \((C,B)\) and push a new
waiting frame whose claim is \(H\).
\item \emph{Defend by \(H\)}: \(H\in |L|\) and \(H D_L^s B\).
The active claim frame closes.
\item \emph{Neutralize using \(H\)}: \(H\in |L|\) and
\(\neg(B D_L C)\). The active claim frame closes.
Here \(H\) names licensed support being cited; the comparison
uses the whole licensed list, not only \(H\)'s conclusions.
\end{enumerate}
On closing a frame, append its claim to \(L\) and remove the frame.
If there is a suspended parent, resume its same open challenge.
Closing a lemma hearing does not close that parent automatically.
If the last frame closes, the proponent wins.
If a proponent position has no legal move, the proponent loses.
Every infinite play also counts as a proponent loss.
The opponent cannot pass when a challenge exists.

Fresh hearings never repeat a licensed or pending claim.
Already licensed arguments may be cited repeatedly.
Different hearings may share subarguments, rule heads or rules;
every individual argument must still satisfy the source grammar.
An opponent can repeat an argument in different frames.
\end{hbdefinition}

The changes from the original dialogue are now explicit. Opponent
challenges still use \(D_\varnothing\). A proponent may suspend a
challenge to establish an auxiliary argument, then return to the
same challenge. Resolutions use already licensed conclusions through
\(D_L\), replacing the original response's self-context
\(D_{\{H\}}\). Global minimality is removed. Prohibition of all
proponent repetitions is replaced by fresh-hearing nonrepetition
and permitted reuse of licensed support.
The argument language, strict rules, generated universe, attack,
dynamic operator and semantic justified set are unchanged.

\begin{hbdefinition}{def:p18-is-licensing-strategy}{Winning strategies and finite proofs}
A proponent strategy selects a legal move at each of its reachable
positions. Its strategy tree contains every legal opponent choice.
It is winning when every resulting maximal play is finite and
ends with the last frame closed. Such a tree is well-founded:
it has no infinite branch. It need not have finitely many nodes
when the theory has infinitely many attackers.
A \emph{finite licensing proof} is a winning strategy tree with
finitely many nodes.
An opponent winning strategy forces a proponent loss under the
same finite-or-infinite outcome convention.
\end{hbdefinition}

\begin{hbtheorem}{thm:p18-is-licensing-sound}{proved=annex:p18.is-licensing-sound}{General soundness and an opponent strategy}
For every source dynamic theory on a set of arguments, if the
proponent has a winning strategy in the corrected dialogue from
empty \(L\) for \(A_0\), then \(A_0\in J\).
If \(A_0\notin J\), the opponent has a winning strategy.
This result does not assume that every arbitrary play has a
semantically sound licensed list.
\end{hbtheorem}

\begin{hbtheorem}{thm:p18-is-licensing-complete}{proved=annex:p18.is-licensing-complete}{Completeness on the full source domain}
For every source dynamic theory,
\[
 A_0\in J
 \quad\Longleftrightarrow\quad
 P_{\rm pl}\text{ has a winning licensing strategy for }A_0.
\]
The proponent strategy can maintain \(|L|\subseteq J\) and choose
fresh hearings with strictly smaller semantic entry rank than
the active claim. It therefore handles shared subarguments and
alternative head producers without any support-compression lemma.
Every source subargument of a claim with a winning strategy also
has a winning strategy.
The two players' outcomes are determined under the specified
infinite-play convention.
\end{hbtheorem}

\begin{hbtheorem}{thm:p18-is-licensing-finite}{proved=annex:p18.is-licensing-finite}{Finite, finitary, and general strategy scope}
For a finitary theory, \(A_0\in J\) if and only if \(A_0\) has
a finite licensing proof.
For a finite argument universe every play is finite, regardless
of the players' choices. For a finite set of ground rules the
winner and a winning strategy can be computed by finite search.
For arbitrary theories the preceding completeness theorem concerns
well-founded strategies, not necessarily finite proof trees.
Finitariness alone is not an effective presentation of the entire
argument universe.
\end{hbtheorem}

\section{A complete repaired exchange}

\begin{hbexample}{def:p18-is-licensing-counterdialogue}{Resolving the formerly impossible challenge}
Use the countertheory and root \(A=[t,z]\). Suppose the opponent
chooses \(B=[m,n]\). The following sequence establishes the needed
support separately:
\begin{enumerate}
\item The proponent opens a hearing for \([]\).
It has no opponent challenge, is licensed, and control returns
to \((A,B)\).
\item The proponent opens a hearing for \(Q=[v,q]\).
It too is unattacked, so it is licensed and the parent resumes.
\item The proponent opens a hearing for \(P=[u,p]\).
Every possible opponent challenge to this hearing is incoherent:
it must derive \(\neg a\) from arguments containing both \(a,b\).
The proponent cites the already licensed empty argument as a
strict defender. The hearing closes and licenses \(P\).
\item Both priorities are now available in \(E_L\).
They neutralize \(B\)'s two rebuttals against \(A\).
The proponent cites licensed \(P\) with the neutralization tag,
closes the root frame, and wins.
\end{enumerate}
No argument containing the inconsistent union of \(P,Q\) has
been used. The only opponent variation inside this exchange is
the choice of incoherent challenge to the \(P\) hearing; every
such choice receives the declared empty-argument defense.
For any different initial challenge to \(A\), the same licensed
support either neutralizes its defeasible negative conclusion or
strictly defeats its incoherent strict-consequence argument.
Thus the exchange extends to a full winning strategy, not merely
one successful branch.
\end{hbexample}

\begin{hbtheorem}{thm:p18-is-licensing-worked}{proved=annex:p18.is-licensing-worked}{The repaired exchange and reuse}
The described strategy wins against every generated opponent
argument in the countertheory.
It remains legal if a hearing for \(U=[u]\) is completed before
the hearing for \(P=[u,p]\); those two licensed arguments share
the same rule and conclusion.
The optional alternative \(u'\) likewise introduces no need to
combine incompatible producers.
The following single-play example has its stated justified set
and completed transcript, but its root has no winning strategy.
\end{hbtheorem}

\begin{hbexample}{def:p18-is-licensing-single-play}{A completed play is not a winning strategy}
Take no ordinary strict rules or priorities and just
\[
 r_a:\Rightarrow a,\quad r_{\neg a}:\Rightarrow\neg a,\quad
 r_b:\Rightarrow b,\quad r_s:\mathord{\sim}b\Rightarrow\neg a.
\]
Here \(J=\{[],[r_b]\}\).
If the opponent challenges \([r_a]\) using \([r_s]\),
the proponent can license \([r_b]\), cite its strict undercut
of \([r_s]\), and close the root in that play.
But \([r_a]\notin J\), and this is not a winning strategy:
the opponent can instead choose \([r_{\neg a}]\).
The complete soundness theorem quantifies over every opponent
choice. A cooperative transcript does not establish its conclusion.
\end{hbexample}

\section{A dynamic game for overruling and defensibility}

\begin{hbdefinition}{def:p18-is-licensing-reversed}{The overruling challenge protocol}
Fix a target argument \(A_0\). A challenger first chooses a legal
argument \(B\) that attacks \(A_0\), using the source attack
relation, which is independent of priority support.
If no such argument exists, the challenger loses.
The challenger then acts as proponent in the corrected licensing
dialogue for \(B\), starting with empty licensed list, while the
defender acts as opponent.
The challenger wins exactly when that dialogue closes its root.
The defender wins when the challenger is unable to continue or
the resulting play is infinite.
This protocol reverses the task of establishing an attacking
argument. It is a newly specified dynamic protocol, not an
unstated exchange of the original dialogue's move rules.
\end{hbdefinition}

\begin{hbtheorem}{thm:p18-is-licensing-reversed}{proved=annex:p18.is-licensing-reversed}{Exact attack-based status correspondence}
The challenger has a winning strategy precisely when \(A_0\)
is overruled in the source's dynamic classification:
\[
 \exists B\in J\quad B\text{ attacks }A_0.
\]
Otherwise the defender has a winning strategy.
Consequently \(A_0\) is source-defensible precisely when the
defender wins this protocol and \(A_0\notin J\).
A justified \(A_0\) also gives a defender win, so that win alone
does not distinguish justified from residual defensible arguments.

For a finitary theory a successful challenger can supply a finite
licensing proof of its attack witness.
No finite certificate for every defender win is asserted on an
infinite universe. On a finite ground-rule domain the whole
classification is decidable by the specified finite games.
\end{hbtheorem}

The distinction between attack and defeat is necessary. An
incoherent argument can receive the designated defeat from
the empty argument without being attacked by it. An initial
challenge in the last protocol must really be an attack.
In the countertheory, \(A\) is justified and its defender wins;
\(B=[m,n]\) is overruled because the justified \(A\) attacks it.
For the earlier two-default inconsistent combination, the empty
argument does not qualify as an attack witness, matching the
source's residual classification rather than a stronger claim
of rational acceptability.

\section{What the correction establishes}

The original finite completeness claim is refuted on its full
strict-rule domain. Therefore it has no unrestricted finitary
extension in that form. The corrected sequential game is sound
and complete on the full set-sized theory domain, using
well-founded strategies, and has finite winning proof trees
under the source's finitary hypothesis.
Every argument remains an actual source argument; only the
dialogue protocol changes. The proof never treats an incompatible
union, a family of argument trees, or a certificate node as one
legal response argument.
Questions about the separate priority-sensitive direct-exclusion
semantics and the historical implementation's computability
scope retain the boundaries recorded in the preceding chapter.
